% !PRINTED 174
\weberpage{174}
We now take as domain of rationality $\Omega$ the field which
consists of the absolute invariants of the group $S$ and all
numbers\footnote{One may also restrict oneself to a particular number field, provided only that the substitution coefficients of the group $S$ are contained in it.}. To this domain of rationality belong the coefficients
of $\Phi(t)$, and we first prove the theorem:

\gesperrt{3. The function $\Phi(t)$ is irreducible in $\Omega$.}

For if $\Psi(t)$ is any function in $\Omega$ which vanishes for $t = \Theta$,
then in the equation $\Psi(\Theta) = 0$, since
the $x_1, x_2, \dots, x_n$ are independent variables, we may replace the system $(x)$
of these variables by $A(x)$, where $A$ is any substitution
of $S$. In this way $\Theta$ can be carried over into each of the functions
$\Theta_1, \Theta_2, \dots, \Theta_{\mu-1}$, while the coefficients
of $\Psi$ remain unchanged, and consequently $\Psi(\Theta_1) = 0$,
$\Psi(\Theta_2) = 0, \dots \Psi(\Theta_{\mu-1}) = 0$. Hence $\Psi(t)$ must be divisible by $\Phi(t)$,
whereby the irreducibility is proved.

Now let $\omega$ denote any function of the variables $x$, and let
\[ \omega,\ \omega_1,\ \omega_2,\ \dots,\ \omega_{\mu-1} \]
be the functions which arise from $\omega$ by application of the
substitutions of $S$; these functions need not all be different from one another.
If $t$ is a variable, then
\weq{7}{\Phi(t)\left(\frac{\omega}{t-\Theta} + \frac{\omega_1}{t-\Theta_1} + \cdots + \frac{\omega_{\mu-1}}{t-\Theta_{\mu-1}}\right) = \Psi(t)}
is an integral rational function of degree $(\mu-1)$ in $t$, and at the same time
it is an invariant of $S$, hence a function in $\Omega$.
If we set $t = \Theta$ in it, then by an inference already often
employed before (Bd.\ I, \S.\,143, 155) it follows that:
\weq{8}{\omega = \frac{\Psi(\Theta)}{\Phi'(\Theta)},}
in which the following theorem is contained:

\gesperrt{4. Every rational function of the variables $(x)$ can
be expressed rationally in terms of $\Theta$, and hence
belongs to the field $\Omega(\Theta)$.}

From this theorem we can derive a second proof
that all invariants of the group $S$ can be represented rationally by a
finite number of them, namely the coefficients of the function
$\Phi(t)$. For if we take an absolute in-
%CONT>>
